\section*{अध्याय \ref{ch:g10:chemical-species} --- रासायनिक स्पीशीज़, प्राकृतिक और संश्लिष्ट}

\begin{solution}{exo:g10:chemical-species:1}
प्राकृतिक: फली से मिला वैनिलिन, कॉफ़ी के बीजों से मिला कैफ़ीन। संश्लिष्ट:
ऐस्पिरिन, कारख़ाने का वैनिलिन।
\end{solution}

\begin{solution}{exo:g10:chemical-species:2}
वे एक ही \omterm{def:g6:pure-substances-and-mixtures:species}{रासायनिक स्पीशीज़} हैं, और किसी स्पीशीज़ के गुण उसके स्रोत से
स्वतंत्र होते हैं।
\end{solution}

\begin{solution}{exo:g10:chemical-species:3}
$R_f = 2.4/6.0 = 0.40$।
\end{solution}

\begin{solution}{exo:g10:chemical-species:4}
वह वाष्प को ठंडा करके फिर द्रव बना देता है। निचले सिरे से प्रवेश करने पर पानी
पूरे आवरण को भर देता है (और सबसे गरम वाष्प से सबसे आख़िर में मिलता है, जिससे
सबसे अच्छी ठंडक मिलती है)।
\end{solution}

\begin{solution}{exo:g10:chemical-species:5}
स्पीशीज़ को पानी से कहीं बेहतर घोले; पानी के साथ \omterm{def:g6:solutions-and-solubility:miscible}{अमिश्रणीय} हो;
जितना संभव हो कम ख़तरनाक हो और आसानी से हटाया जा सके।
\end{solution}

\begin{solution}{exo:g10:chemical-species:6}
एथेनॉल पानी के साथ \omterm{def:g6:solutions-and-solubility:miscible}{मिश्रणीय} है: कोई परत नहीं बनती, कुछ अलग नहीं किया जा सकता।
साइक्लोहेक्सेन पानी के साथ \omterm{def:g6:solutions-and-solubility:miscible}{अमिश्रणीय} है और लिनालूल को अच्छी तरह घोलता है।
\end{solution}

\begin{solution}{exo:g10:chemical-species:7}
पानी ऊपर है: डाइक्लोरोमेथेन के \qty{1}{mL} का भार \qty{1.33}{g} है,
जो \qty{1}{mL} पानी से अधिक है, इसलिए वह नीचे बैठता है।
\end{solution}

\begin{solution}{exo:g10:chemical-species:8}
वह बहुत कम ताप पर (ऐस्पिरिन: \qty{135}{\celsius}) और एक परास में पिघलता है:
नमूना अशुद्ध है, या ऐस्पिरिन है ही नहीं।
\end{solution}

\begin{solution}{exo:g10:chemical-species:9}
लिनालूल और लिनालिल एथेनोएट (उसके दोनों धब्बे दोनों संदर्भों की ऊँचाइयों
पर हैं)। लिनालिल एथेनोएट ज़्यादा ऊपर चढ़ा।
\end{solution}

\begin{solution}{exo:g10:chemical-species:10}
\omterm{def:g10:chemical-species:distillation}{जल-आसवन}; आसुत में साइक्लोहेक्सेन मिलाकर हिलाना;
परतों को अलग करना; साइक्लोहेक्सेन का \omterm{def:g3:separating-mixtures:evaporate}{वाष्पन}।
\end{solution}

\begin{solution}{exo:g10:chemical-species:11}
पेंसिल \omterm{def:g10:chemical-species:tlc}{निक्षालक} में नहीं \omterm{def:g2:mixing-and-dissolving:dissolve}{घुलती}; \omterm{def:g10:chemical-species:tlc}{निक्षालक} से ऊपर होने पर धब्बे टंकी में \omterm{def:g2:mixing-and-dissolving:dissolve}{घुलने} के
बजाय प्लेट पर ऊपर ले जाए जाते हैं।
\end{solution}

\begin{solution}{exo:g10:chemical-species:12}
साइक्लोहेक्सेन \qty{81}{\celsius} पर उबलता है, लिनालूल \qty{198}{\celsius} पर:
हल्का गरम करने से साइक्लोहेक्सेन निकल जाता है जबकि लिनालूल बचा रहता है।
\end{solution}

\begin{solution}{exo:g10:chemical-species:13}
नमूना शायद वही स्पीशीज़ है (उसी प्लेट पर \omterm{def:g10:chemical-species:natural}{प्राकृतिक स्पीशीज़} जितना $R_f$) और शायद
शुद्ध है (एक धब्बा, तीक्ष्ण गलनांक)।
\end{solution}

\begin{solution}{exo:g10:chemical-species:14}
$0.40 \times 5.5 = \qty{2.2}{cm}$। दूसरा $0.75 \times 5.5 =
\qty{4.1}{cm}$ पर (निकटतम mm तक), यानी A से लगभग \qty{1.9}{cm} ऊपर।
\end{solution}

\begin{solution}{exo:g10:chemical-species:15}
\qty{200}{mL} पानी $1.6 \times 0.200 = \qty{0.32}{g}$ घोल सकता है:
\qty{0.20}{g} के साथ वह संतृप्त नहीं है; वह \qty{0.12}{g} और घोल सकता है।
घुला हुआ \qty{0.20}{g} पानी के साथ खो जाता; साइक्लोहेक्सेन,
जो लिनालूल को कहीं बेहतर घोलता है, उसका ज़्यादातर भाग वापस ले लेता है।
\end{solution}

\begin{solution}{pb:g10:chemical-species:1}
\textbf{1.} उसकी भाप सुगंधित स्पीशीज़ को फूलों से निकालकर संघनित्र में
ले जाती है।

\textbf{2.} वह वाष्प को ठंडा करके फिर द्रव बनाता है; ठंडा पानी निचले सिरे से
प्रवेश करता है।

\textbf{3.} \omterm{def:g6:solutions-and-solubility:miscible}{अमिश्रणीय}; पानी से कम घना (वह तैरता है)।

\textbf{4.} परत बहुत पतली है और तेल का कुछ भाग पानी में घुला या बिखरा हुआ है:
\omterm{def:g10:chemical-species:extraction}{निष्कर्षण} उसे वापस पा लेता है।

\textbf{5.} वह लिनालूल को अच्छी तरह घोलता है, पानी के साथ \omterm{def:g6:solutions-and-solubility:miscible}{अमिश्रणीय} है, और कम ताप
(\qty{81}{\celsius}) पर उबलता है, इसलिए उसे हटाना आसान है।

\textbf{6.} नहीं: एथेनॉल पानी में मिल जाता है और कोई परत नहीं बनती।

\textbf{7.} ऊपर (\qty{0.78}{g} प्रति mL, पानी से कम)।
डाइक्लोरोमेथेन (\qty{1.33}{g} प्रति mL) के साथ \omterm{def:g6:solutions-and-solubility:solute}{विलायक} की परत
नीचे होती।

\textbf{8.} पास में कोई लौ नहीं (ज्वलनशील); धूम-हुड के नीचे दस्ताने पहनकर
काम करना, और अपशिष्ट को नाली में बहाने के बजाय इकट्ठा करना
(जलीय जीवों के लिए विषैला)।

\textbf{9.} साइक्लोहेक्सेन \qty{81}{\celsius} पर उबलता है, लिनालूल
\qty{198}{\celsius} पर: हल्का गरम करने से केवल \omterm{def:g6:solutions-and-solubility:solute}{विलायक} हटता है।

\textbf{10.} $R_f(\text{L}) = 2.4/6.0 = 0.40$; $R_f(\text{A}) = 4.5/6.0
= 0.75$।

\textbf{11.} लिनालूल और लिनालिल एथेनोएट।

\textbf{12.} नहीं: उसमें कम से कम दो स्पीशीज़ हैं; वह एक \omterm{def:g6:pure-substances-and-mixtures:mixture}{मिश्रण} है।

\textbf{13.} ताकि सभी धब्बे एक जैसी परिस्थितियों में विकसित हों:
तभी ऊँचाइयों की तुलना की जा सकती है।

\textbf{14.} उसका $R_f$ उन्हीं परिस्थितियों में लिनालूल के बराबर है: वह
शायद लिनालूल है, \omterm{def:g10:chemical-species:natural}{प्राकृतिक स्पीशीज़} के समान।

\textbf{15.} $1.0 \times 0.90 = \qty{0.90}{g}$।

\textbf{16.} $0.90/100 = \qty{0.90}{\percent}$।

\textbf{17.} $100/0.009 \approx 11\,000$\,g, यानी लगभग \qty{11}{kg}
फूल।

\textbf{18.} वह बिना \omterm{def:g5:irreversible-changes:chemical-change}{रासायनिक परिवर्तन} के एक पौधे से निकाला गया है। पर उसका
लिनालूल संश्लिष्ट लिनालूल वाली ही स्पीशीज़ है: वही सूत्र, वही
गुण।

\textbf{19.} \qty{100}{g} फूलों से \qty{0.90}{g} सगंध तेल।
\end{solution}
